% =============================================================================
%  CHAPTER 24 -- PHASE 11: ARCHITECTURE AND LATENCY
% =============================================================================
\section{Phase 11 --- System architecture and latency analysis}
\label{ch:phase11}

\textbf{Reproduce it:} \code{cmake --build build -j2}, \code{./build/statarbsim --demo},
\code{./build/microbench}, then \code{python3
scripts/analysis/capture\_engine\_bench.py \&\& python3
scripts/analysis/make\_book\_tables.py \&\& python3
scripts/plotting/render\_engine\_results.py}.\\
\textbf{Math used:} Chapter~\ref{ch:systems}.\\
\textbf{Artifacts:} \code{src/} (all), \code{bench/microbench.cpp},
\code{results/tables/engine\_\{bench,latency\}.csv},
\code{results/figures/\{latency\_hist,cache\_bench\}.png},
\code{report/tables/engine\_\{bench,latency\}.tex}.

\subsection{Goal}

Phase~\ref{ch:phase3} built the engine and measured it once. This phase makes the systems
side of the project a first-class result: write down the design rationale with the
cache-line arithmetic behind each choice, re-measure every claim from a live run, make each
number traceable to a capturing script, and report the multipliers as measured rather than
as the theory would like.

\subsection{Architecture as shipped}
\label{sec:arch:principles}

The data flow is a straight causal pipeline, one module per stage, no allocation after
construction:

\begin{center}
\small
\begin{verbatim}
CSV (data/processed/engine/<T>.csv)
   |  src/data/csv_loader.h
   v
MarketDataBuffer  (SoA ring: ts[], px[], vol[] contiguous)     src/data/market_data_buffer.h
   |  push bar t-1
   v
RollingRegression (rank-1 XtX / Xty accumulators, flat row ring) src/model/rolling_regression.h
   |  Cholesky (+ lambda I for ridge)                            src/model/dense_la.h
   |  [alternative estimators: BasketSelector OLS/Ridge/Lasso/EN/PCR, Jacobi eigen]
   v
SignalGenerator   (running mean/var -> z, two-band hysteresis)   src/signal/signal_generator.h
   |  desired direction for bar t+1
   v
PortfolioBook     (positions, gross/net exposure, mark-to-market) src/portfolio/portfolio_book.h
   |  [multi-book layer: risk.h inverse-vol weights + limits]
   v
LatencyHist x3    (regression_update, signal_generation, order_placement)
                  src/engine/latency_hist.h
\end{verbatim}
\end{center}

Each design principle is a claim about memory, and each is testable:

\begin{itemize}
  \item \textbf{Contiguous flat storage, never pointer-chasing containers.} A cache miss
    fetches a 64-byte line; with contiguous data the line contains 8 useful doubles and the
    prefetcher streams ahead. \code{std::map}/\code{std::list} would stride through
    scattered heap lines. \emph{Test:} \code{buffer\_basic\_and\_contiguous}.
  \item \textbf{Structure-of-Arrays.} Timestamps, prices and volumes live in separate
    contiguous buffers, so a loop that needs only prices touches one column
    (Section~\ref{sec:sys:soa}). \emph{Measured:} SoA vs AoS below.
  \item \textbf{No dynamic allocation on the hot path.} All buffers sized in the
    constructor; \code{update()} and \code{compute()} write into pre-allocated scratch.
    \emph{Test:} \code{no\_allocation\_in\_hot\_push\_path} --- a global
    \code{operator new} counter reads 0 after 1{,}000 windowed updates and solves.
  \item \textbf{Incremental regression, not refits.} Rank-1 accumulators
    \eqref{eq:rank1update}: $O(n^2)$ per bar instead of $O(Wn^2)$, with a single small
    Cholesky solve when coefficients are requested. \emph{Test:} matches brute force to
    $10^{-8}$ at every window position, and the eviction test proves the right row leaves.
  \item \textbf{No virtual dispatch on the hot path.} Per-bar units are concrete,
    header-only, template-friendly classes; an indirect call would cost 2--5\,ns and block
    inlining --- 10--25\% of the cheapest stage.
  \item \textbf{Dependency-free linear algebra.} Cholesky + cyclic Jacobi in
    \code{dense\_la.h}, chosen because Eigen is not installable and not vendable at a
    reviewable size. \emph{Tests:} $LL^{\top}=A$, ridge inverts a singular system,
    $V^{\top}AV=\Lambda$, $\mathrm{PCR}(k{=}n)\equiv\mathrm{OLS}$.
  \item \textbf{Risk primitives as pure functions.} \code{src/portfolio/risk.h} provides
    inverse-vol weights \eqref{eq:inversevol} and gross/net/concentration limit
    verification, allocation-free and unit-tested independently of market data, and is
    consumed by the Phase~\ref{ch:phase7} book.
\end{itemize}

\subsection{Measured cache microbenchmarks}

\input{tables/engine_bench.tex}

\begin{figure}[H]\centering
\includegraphics[width=0.92\linewidth]{figures/cache_bench.png}
\caption{Cache-efficiency microbenchmarks, measured from \code{build/microbench}: the
rolling-window sweep over a contiguous \code{std::vector} versus a \code{std::deque} at a
250-bar window, and a Structure-of-Arrays versus Array-of-Structures column reduction over
$2^{20}$ elements.}
\label{fig:cache_bench}
\end{figure}

\begin{itemize}
  \item \textbf{\code{vector} vs \code{deque}} (rolling-window sweep, 250-bar window): the
    contiguous \code{vector} wins in every build measured, by a ratio that has landed
    between $\mathbf{1.08\times}$ and $\mathbf{1.33\times}$ across sessions
    ($1.084\times$, $1.264\times$, $1.333\times$ in three consecutive builds of the same
    binary on the same data). The live value is in the table above; the honest summary is
    ``measurably faster, not reliably by how much''.
  \item \textbf{SoA vs AoS} (column reduction over $2^{20}$ elements): likewise always in
    favour of SoA, with a ratio observed between $\mathbf{1.03\times}$ and
    $\mathbf{1.51\times}$ across builds. The absolute cost of the reduction
    ($\approx1.5$--$2.2$\,ms per pass over 8\,MB / 24\,MB working sets) is dominated by
    memory traffic, which is exactly why it is the noisiest number in the repository.
\end{itemize}

\begin{remark}[Honesty: the multipliers are smaller than theory, and unstable]
Sections~\ref{sec:sys:soa} and \ref{sec:sys:vector} derive upper bounds of $\approx3\times$
for both comparisons (three times fewer line fetches for SoA; double indirection and broken
stride for \code{deque}). The measurements fall short of that bound and, worse, are not
reproducible to better than tens of percent. Three concrete reasons, in order of size:
\begin{enumerate}
  \item \textbf{Both working sets partly fit in cache.} At a 250-bar window the
    \code{deque} blocks and the \code{vector} are both L1/L2-resident, so the penalty is a
    few cycles of address arithmetic per access, not a DRAM miss. The theory's $3\times$
    assumes memory-bound streaming, which this microbenchmark does not create.
  \item \textbf{We do not control the machine.} This is a 2-core shared sandbox with its
    own neighbours on the same L3 and the same memory controller; a run started while
    another process is compiling gives a different answer. The spread across builds
    ($1.03$--$1.51\times$ for SoA/AoS) is larger than the effect being measured in the
    favourable direction, so the only defensible claim is the sign.
  \item \textbf{The benefit is architectural, not a single multiplier.} SoA and contiguous
    storage pay off when the book spans hundreds of baskets (working sets that do
    \emph{not} fit in cache), when several operations stream the same column, and when the
    compiler can vectorise --- none of which a two-line microbenchmark exercises.
\end{enumerate}
The tables state whatever the last run measured, and the prose states the range across runs.
Neither is rounded up to ``2--3$\times$, as theory suggests''.
\end{remark}

\subsection{Measured per-stage latency}

\input{tables/engine_latency.tex}

\begin{figure}[H]\centering
\includegraphics[width=0.82\linewidth]{figures/latency_hist.png}
\caption{Per-stage latency distribution (p50/p95/p99, ns) of the C++ pipeline over
30{,}000 synthetic bars, from \code{build/statarbsim --demo}.}
\label{fig:latency_hist}
\end{figure}

Reading of Table~\ref{tab:engine_latency}, including its limitations:
\begin{itemize}
  \item \textbf{The regression stage dominates}: p50 $\approx0.6\,\mu$s versus
    $\approx21$\,ns for signal generation and order placement --- a ratio of
    $\approx25$--$30\times$, stable across builds even though the absolute values are not
    (an earlier build recorded p50 $\approx518$\,ns and $\approx28$\,ns for the cheap
    stages). That is where any optimisation effort belongs, and the ratio is consistent
    with the flop count of Section~\ref{sec:sys:complexity} (the Cholesky plus two rank-1
    updates plus bookkeeping).
  \item \textbf{The quantiles are histogram bucket edges, not raw order statistics.} The
    instrument has 64 logarithmically spaced buckets over
    $[10\,\text{ns},100\,\mu\text{s}]$, so each bucket spans a factor
    $10^{4/63}=1.158$, and the reported quantile is the \emph{lower edge} of the bucket
    containing it. $599.5$\,ns is such an edge (the next is $694$\,ns, then $804$\,ns,
    $931$\,ns, $1078$\,ns). Two consequences follow, and both have been observed: when the
    distribution is tight, p50, p95 and p99 all report the \emph{same} edge (an earlier
    build printed $599.5$ for all three, which is a resolution artifact, not a degenerate
    distribution); when it is not, they step up the ladder of edges, as in the build that
    produced the table above ($599.5 / 1075.8 / 2234.6$\,ns). The mean, which is computed
    from the raw samples and not from the histogram, lies above the p50 edge in every build
    ($660$--$766$\,ns observed) --- confirming mass spread through and beyond it. Reporting
    bucket edges rather than inventing intermediate values is the honest convention;
    tightening the range to $[100\,\text{ns},10\,\mu\text{s}]$ would resolve the tail and is
    listed as a small improvement in Chapter~\ref{ch:conclusion}.
  \item \textbf{The two cheap stages are at the measurement floor.} \code{chrono} reads cost
    $\approx$20--25\,ns via the vDSO, and both stages report $\approx20.8$\,ns --- i.e.\ we
    are measuring the instrument as much as the work. This is disclosed rather than
    presented as ``21\,ns signal generation''.
\end{itemize}

\subsection{Traceability: how every systems number reaches the page}

\begin{center}
\renewcommand{\arraystretch}{1.2}
\small
\begin{tabular}{p{0.30\linewidth}p{0.32\linewidth}p{0.30\linewidth}}
\toprule
Number & Produced by & Lands in \\
\midrule
p50/p95/p99 per stage, mean & \code{./build/statarbsim --demo} $\to$ \code{scripts/analysis/capture\_engine\_bench.py} & \code{results/tables/engine\_latency.csv} $\to$ \code{report/tables/engine\_latency.tex}, \code{figures/latency\_hist.png} \\
\code{vector} vs \code{deque}, SoA vs AoS & \code{./build/microbench} $\to$ same capturing script & \code{results/tables/engine\_bench.csv} $\to$ \code{report/tables/engine\_bench.tex}, \code{figures/cache\_bench.png} \\
Test counts and assertions & \code{./build/statarbsim --test} & \code{environment/check.sh} gate \\
Allocation-free property & \code{test\_no\_allocation\_in\_hot\_push\_path} & same \\
\bottomrule
\end{tabular}
\end{center}

Nothing in this chapter's tables is typed by hand, and the whole set is re-measured on
every \code{python3 report/build\_report.py}. That has a visible consequence: the numbers
move between runs, sometimes a lot (the regression p50 has been measured at $518$\,ns and
at $599.5$\,ns; the \code{deque} ratio at $1.084\times$, $1.264\times$ and $1.333\times$;
the SoA/AoS ratio at $1.033\times$ and $1.509\times$). The prose above therefore quotes
ranges and mechanisms, and leaves the point estimates to the generated tables, which always
show the most recent run. Reporting the live values rather than the remembered ones is the
point of the rule --- even when, as here, the live values are noisier and less flattering
than an earlier session's.

\subsection{Exit gate --- status}

\begin{itemize}
  \item[\checkmark] \code{vector} vs \code{deque} and SoA vs AoS microbenchmarks generated
    from real runs;
  \item[\checkmark] latency histograms (p50/p95/p99) plotted from a live
    \code{--demo} run, with the resolution limitation stated;
  \item[\checkmark] design rationale written with the cache-line arithmetic, not asserted;
  \item[\checkmark] every benchmark number traceable to a capturing script and a CSV;
  \item[\checkmark] the C++ test suite (24 tests, 2{,}852 assertions) and \code{check.sh}
    (52/52) green.
\end{itemize}
